Combining many antenna elements lets a system synthesize and steer a beam
electronically --- the basis of modern radar and 5G.

\section{The Array Factor}
For $N$ identical elements spaced $d$ along a line, the far-field pattern is the
element pattern times the array factor
\[ \text{AF}(\theta) = \sum_{n=0}^{N-1} e^{\,j n (k d \cos\theta + \beta)}, \]
where $\beta$ is the progressive phase between elements. The array factor, not the
element, gives the array its directivity.

\section{Steering}
Choosing $\beta = -k d \sin\theta_0$ (for a broadside reference) points the main
beam to angle $\theta_0$ --- with no moving parts, just phase. The beam can be
repointed in microseconds, which is why phased arrays replaced mechanically
scanned dishes in radar.

\section{Beamwidth, Grating Lobes, Tapering}
More elements and wider aperture narrow the beam (beamwidth $\propto \lambda/(Nd)$).
If $d>\lambda$, extra \emph{grating lobes} appear (usually $d\le\lambda/2$ avoids
them). Amplitude tapering across the array trades a little beamwidth for much lower
sidelobes.

\section{Digital Beamforming and MIMO}
With a receiver per element, beams are formed in DSP --- many simultaneous beams,
adaptive nulling of interferers, and, in communications, spatial multiplexing
(MIMO/massive MIMO) that sends independent streams on the same frequency.
